\section{What each input contributes}
\label{sec:inputs}

This section asks a different question from the rest of the report: not how well
a policy performs, but which of its inputs it is actually using.
Figure~\ref{fig:streams} carries the answer. The method, its one essential
precaution and the rule for reading the figure come first, because the figure
is easy to over-read.

Each input stream is suppressed in turn --- replaced by an uninformative
substitute --- and the resulting change in the planned action is measured. A
stream's share is therefore the size of the difference that stream makes to
this trained policy on these frames. For a policy whose output is sampled
rather than determined, the sampler must be fixed while this is done, and the
reason is not fastidiousness: on the diffusion family, two plans drawn from one
unchanged observation differ by more than eighty times the difference that
suppressing a camera makes. Unpinned, the measurement would report the
sampler's own variation and call it attribution.

\begin{figure}[t]
  \centering
  \includegraphics[width=\textwidth]{stream_shares.png}
  \caption{The share of the change in the planned action attributable to each
  group of inputs. Shares are normalised within a policy, so each bar sums to
  one hundred and the \emph{heights are not comparable between bars}. What may
  be carried across is the order of the groups within each bar. The third
  family reads two fingertip cameras where the others read four, because it
  finetunes a pretrained model with a fixed number of image slots.}
  \label{fig:streams}
\end{figure}

The rule for reading Figure~\ref{fig:streams} is that only the \emph{ordering}
within a bar may be carried between bars. The shares are normalised per policy,
so a larger slice in one bar than in another says nothing on its own; what it
can say is which input a given policy leans on most, and where the tactile
channels sit in its ranking.

By that rule the families disagree, and not subtly. The action-chunking family
leans on the overhead camera, with proprioception second and the four tactile
cameras last by a wide margin. The diffusion family leans overwhelmingly on
proprioception --- where the arms are matters more to it than what any camera
shows --- with vision and touch far behind. The third family ranks touch
\emph{above} proprioception, the only one of the three to do so, and that is an
ordering claim the normalisation permits. It also reads half as many tactile
cameras as the others, so whether that ranking reflects the architecture or the
smaller, differently-chosen input set cannot be told apart here.

Two limits bound all of this. The frames are drawn from recordings the policies
trained on, so the figure describes where a trained policy looks on familiar
material rather than on the held-out recordings the rest of the report scores.
And a share measures how much suppressing an input moves the plan, which is not
the same as how much the input matters for folding a garment: a stream can be
ignored because another carries the same information, and only retraining
without it would show whether it was needed. Section~\ref{sec:threats} returns
to both.
